\documentclass[12pt]{article}
\usepackage[utf8]{inputenc}
\usepackage[T1]{fontenc}
\usepackage{hyperref}
\title{Mitlatongo Mixtec Language Resources for K-12 Curriculum Development}
\author{Compiled by Claude Sonnet 5 (claude-sonnet-5) for Andr\'as Kornai}
\date{}
\begin{document}
\maketitle

\section{Introduction}
Mitlatongo Mixtec (ISO 639-3 vmm; Glottocode mitl1235) is an Otomanguean language of southern Oaxaca, Mexico, with about 2,800 speakers in the reference description. Searches found only a Wikipedia article and audio pages from the Global Recordings Network. No school, dictionary, or grammar materials were found.

\section{Primary Educational Resources}
None found in search results.

\section{Linguistic and Reference Materials}
\begin{itemize}
\item Wikipedia article on Mitlatongo-Yutanduchi Mixtec. It treats Mitlatongo as one of two varieties, cites about 2,800 speakers, and gives roughly 70\% intelligibility between the two: \url{https://en.wikipedia.org/wiki/Mitlatongo-Yutanduchi_Mixtec}
\end{itemize}

\section{Religious and Community Resources}
No Mitlatongo-specific Bible translation was found. Pages for the nearby Diuxi-Tilantongo Mixtec are a different language and are not used here.

\section{Digital and Audio Resources}
\begin{itemize}
\item Global Recordings Network audio in Mitlatongo Mixtec (evangelistic and Bible-teaching content): \url{https://globalrecordings.net/en/language/vmm}
\end{itemize}

\section{Curriculum Development Context}
Very little documented material is available online. A curriculum would depend on community-based documentation and on checking whether materials exist for the related Yutanduchi variety.
\end{document}